%=========== ΛΥΣΗ - 50ο ΘΕΜΑ Α ===========
\providecommand{\theoriaref}[3]{%
\begin{center}
\fcolorbox{maincolor!70!black}{maincolor!8}{%
\begin{minipage}{.88\linewidth}
\textcolor{maincolor!80!black}{\kerkissans{\textbf{#1~\ref{#2}}}}\\[1mm]
#3
\end{minipage}}
\end{center}}
\providecommand{\orismosref}[1]{%
\theoriaref{Ορισμός}{#1}{Η απάντηση δίνεται από τον αντίστοιχο ορισμό.}}
\providecommand{\apodeixiref}[1]{%
\theoriaref{Θεώρημα}{#1}{Η ζητούμενη απόδειξη δίνεται στην αντίστοιχη απόδειξη.}}

\begin{thema}{Α}
\begin{erwthma}
\item \apodeixiref{thewrhma:synepeies_thmt_2}

\item \orismosref{orismos:katakoryfh_asymptwth}

\item \orismosref{orismos:antistrofh_synarthsh}

\item Οι χαρακτηρισμοί των προτάσεων είναι:
\begin{alist}
\item \textbf{Λάθος}. Για παράδειγμα, η συνάρτηση $f(x)=x^2$ στο $[-1,2]$ είναι συνεχής και παραγωγίσιμη, έχει $f(-1)\neq f(2)$, αλλά $f'(0)=0$.
\item \textbf{Λάθος}. Για παράδειγμα, η $f(x)=1-x$ είναι $1-1$ και είναι ίση με την αντίστροφή της, οπότε οι $C_f$ και $C_{f^{-1}}$ έχουν κοινά σημεία και εκτός της ευθείας $y=x$.
\item \textbf{Σωστό}. Αν η $f$ παρουσίαζε τοπικό ακρότατο στο εσωτερικό σημείο $x_0$ και ήταν παραγωγίσιμη εκεί, από το θεώρημα \eng{Fermat} θα είχαμε $f'(x_0)=0$, άρα το $x_0$ θα ήταν κρίσιμο σημείο.
\item \textbf{Σωστό}. Είναι ο τύπος της ολοκλήρωσης κατά παράγοντες:
\[
\dint_a^\beta f(x)g'(x)\d x
=f(\beta)g(\beta)-f(a)g(a)-\dint_a^\beta f'(x)g(x)\d x.
\]
\item \textbf{Σωστό}. Από το $\lim\limits_{x\to x_0}f(x)=+\infty$, για $M=0$ υπάρχει διάτρητη περιοχή του $x_0$ στην οποία ισχύει $f(x)>0$.
\end{alist}

\item Σωστή απάντηση είναι η \textbf{γ}. Επειδή $f(x)>0$ στο $(0,4)$, το εμβαδόν του χωρίου είναι
\[
\dint_0^4 f(x)\d x=10.
\]
Αν $F$ είναι μια αρχική της $f$, τότε από το θεμελιώδες θεώρημα του ολοκληρωτικού λογισμού έχουμε
\[
F(4)-F(0)=\dint_0^4 f(x)\d x=10.
\]
\end{erwthma}
\end{thema}
%=========== ΤΕΛΟΣ ΛΥΣΗΣ - 50ο ΘΕΜΑ Α ===========
